% GENERATED FILE. Edit canonical lessons, then run npm run generate:books.
% canonical-source-hash: fnv1a64:b39eb39ec60b91ae
% canonical-lessons: GE-C07-tag-muster, GE-C07-mond, GE-C07-montag, GE-C07-dienstag, GE-C07-donner, GE-C07-donnerstag, GE-C07-freitag, GE-C07-mittwoch, GE-C07-wochentage-1

\chapter{The Gods of the Week}
\label{ch:ge-wochentage}
\hlchaptermodality{10}

\begin{chapteropening}
\textbf{By the end of this chapter:} \emph{I can name the five German weekdays, say which god is in each, and explain why one of them has no god at all.}

The Romans named their weekdays for planet-gods and the Germanic peoples swapped in their own, one for one, which is why German and English line up here and Romance does not. Five days are a god's name plus Tag; this chapter meets them one at a time, with the nouns Mond and Donner taught before the days built on them, and ends on Mittwoch — the day the medieval Church took Wodan out of.
\end{chapteropening}

\section[Gott plus Tag]{Gott plus Tag}

\label{lesson:GE-C07-tag-muster}



\begin{quote}
You have now seen a language keep its numbers and buy its months. The weekdays are the same question asked again, and German answers it a third way.
\end{quote}

\begin{grammarlens}[title={Grammar Lens: a god, then Tag}]
\begin{quote}
A German weekday is a \textbf{god's name} with \emph{\textbf{Tag}} on the end.
\end{quote}

\emph{Tag} you already own. What sits in front of it is the part worth knowing: the Romans named their weekdays for planet-gods — Mars, Mercury, Jupiter, Venus — and when those names reached the Germanic peoples, they did not translate them and they did not borrow them. They \textbf{swapped in their own gods}, one for one.

That is why German and English line up here and Romance does not. Both inherited the same substitutions, made once, long before either language was written down.

Five of the seven days keep that shape. Two do not, and both were edited later for reasons you can name — which is the more interesting half of the chapter.
\end{grammarlens}

\subsection*{Your turn}
Say these aloud:

\begin{itemize}
\raggedright
  \item the shape — \textquotedblleft{}a god, then Tag\textquotedblright{}
  \item what the Romans used — \textquotedblleft{}planet-gods\textquotedblright{}
  \item what Germanic did — \textquotedblleft{}swapped in its own, one for one\textquotedblright{}
\end{itemize}

\emph{Twice through:} \textquotedblleft{}Gott plus Tag.\textquotedblright{}

\begin{tcolorbox}[breakable,colback=teal!4,colframe=teal!35!black,title={Before you move on}]
What shape does a German weekday have? (\textbf{A god's name plus \emph{Tag}}.) Whose gods did the Romans use? (\textbf{Their planet-gods}.) Did Germanic borrow those names or replace them? (\textbf{Replaced them}, one for one.) How many of the seven keep the shape? (\textbf{Five}.) Next: the first of them.
\end{tcolorbox}
\section[der Mond]{der Mond — \textquotedblleft{}the moon\textquotedblright{}}

\label{lesson:GE-C07-mond}



\begin{quote}
Before the first weekday, the thing it is named after.
\end{quote}

\subsection*{You'll want to know: der Mond}
\begin{quote}
\textbf{der Mond} — \textquotedblleft{}the moon.\textquotedblright{}
\end{quote}

A \emph{der}-word, and the same inherited word as English \textbf{moon}. German rounded the vowel and added a \emph{d} on the end; nothing else happened to it.

The moon was the first Germanic god-slot in the week, which is why it comes first here.

\begin{sounds}
\begin{itemize}
\raggedright
  \item \texttt{long-o} — a long, round \emph{o}, the vowel of English \emph{moan}, not the \emph{oo} of \emph{moon}.
  \item \texttt{d-t-final} — the final \emph{d} hardens to a \emph{t} at the end of the word, the same hardening you met in \emph{und}. So it comes out \emph{mohnt}.
\end{itemize}
\end{sounds}

\subsection*{Your turn}
Say these aloud:

\begin{itemize}
\raggedright
  \item \textquotedblleft{}der Mond\textquotedblright{} — \emph{mohnt}, with a hard ending
  \item the pair — \textquotedblleft{}Mond, moon\textquotedblright{}
  \item the vowel warning — \textquotedblleft{}not \emph{moon}; \emph{moan}\textquotedblright{}
\end{itemize}

\emph{Twice through:} \textquotedblleft{}der Mond.\textquotedblright{}

\begin{tcolorbox}[breakable,colback=teal!4,colframe=teal!35!black,title={Before you move on}]
Say \textquotedblleft{}the moon.\textquotedblright{} (\textbf{Der Mond}.) Which English word is it? (\emph{\textbf{Moon}}.) How does the final \emph{d} sound? (\textbf{Like a t}.) Next: the day named after it.
\end{tcolorbox}
\section[Montag]{Montag — \textquotedblleft{}Monday\textquotedblright{}}

\label{lesson:GE-C07-montag}



\begin{quote}
You have the moon and you have the day. Put them together before reading on.
\end{quote}

\subsection*{You'll want to know: Montag}
\begin{quote}
\textbf{Montag} — \textquotedblleft{}Monday.\textquotedblright{}
\end{quote}

\emph{Mon-} is \emph{Mond} with its hard ending worn off inside the compound, and \emph{-tag} is the day. The pattern's first instance, and the easiest: English says the same two words in the same order — \emph{moon} and \emph{day}.

That is what an \textbf{inherited} name looks like. Nobody translated this; both languages were already saying it before they were separate languages.

\subsection*{Your turn}
Say these aloud:

\begin{itemize}
\raggedright
  \item \textquotedblleft{}Montag\textquotedblright{}
  \item its two pieces — \textquotedblleft{}Mond, Tag\textquotedblright{}
  \item the English beside it — \textquotedblleft{}Montag, Monday\textquotedblright{}
\end{itemize}

\emph{Twice through:} \textquotedblleft{}Montag — Monday.\textquotedblright{}

\begin{tcolorbox}[breakable,colback=teal!4,colframe=teal!35!black,title={Before you move on}]
Say \textquotedblleft{}Monday.\textquotedblright{} (\textbf{Montag}.) What are its two pieces? (\emph{\textbf{Mond}} \textbf{and} \emph{\textbf{Tag}}.) Is the English name a translation of it? (\textbf{No} — both inherited it.) Next: a day named for a god you have never heard of.
\end{tcolorbox}
\section[Dienstag]{Dienstag — \textquotedblleft{}Tuesday\textquotedblright{}}

\label{lesson:GE-C07-dienstag}



\begin{quote}
The second day, and the first god in the week whose name you will not recognise — though you say it every Tuesday.
\end{quote}

\subsection*{You'll want to know: Dienstag}
\begin{quote}
\textbf{Dienstag} — \textquotedblleft{}Tuesday.\textquotedblright{}
\end{quote}

\emph{Diens-} plus \emph{-tag}, and the \emph{ie} is the long \emph{ee} you met in \emph{vier}: \emph{DEENS-tahk}.

\begin{cousinweb}
The god is \textbf{Tiw}, or \textbf{Ziu} in the older German — the Germanic war-god, put in where the Romans had \textbf{Mars}. English kept his name almost bare: \textbf{Tues}day is \emph{Tiw's day}.

And his name is the oldest thing in this chapter. It comes from the Proto-Indo-European word for the \textbf{bright sky}, and the same word became Latin \emph{deus} (\textquotedblleft{}god\textquotedblright{}), Greek \textbf{Zeus}, and the \emph{Ju-} of \emph{Jupiter}.

So the god in \emph{Tuesday} and the god in \emph{Jupiter} are the same ancient word, which the Germanic peoples had already demoted from sky-king to war-god long before the Romans arrived to be swapped out.
\end{cousinweb}

\subsection*{Your turn}
Say these aloud:

\begin{itemize}
\raggedright
  \item \textquotedblleft{}Dienstag\textquotedblright{} — \emph{DEENS-tahk}
  \item the god — \textquotedblleft{}Tiw, Ziu\textquotedblright{}
  \item the family — \textquotedblleft{}Tiw, deus, Zeus\textquotedblright{}
\end{itemize}

\emph{Twice through:} \textquotedblleft{}Dienstag — Tuesday.\textquotedblright{}

\begin{tcolorbox}[breakable,colback=teal!4,colframe=teal!35!black,title={Before you move on}]
Say \textquotedblleft{}Tuesday.\textquotedblright{} (\textbf{Dienstag}.) Which god is in it? (\textbf{Tiw}, or \emph{Ziu}.) Which Roman god did he replace? (\textbf{Mars}.) Which Greek god shares his name's root? (\textbf{Zeus}.) Next: the noun behind the loudest day of the week.
\end{tcolorbox}
\section[der Donner]{der Donner — \textquotedblleft{}thunder\textquotedblright{}}

\label{lesson:GE-C07-donner}



\begin{quote}
Run the sound law backwards on this one before you read it. German has no \emph{th}; what would English \emph{thunder} look like in German?
\end{quote}

\subsection*{You'll want to know: der Donner}
\begin{quote}
\textbf{der Donner} — \textquotedblleft{}thunder.\textquotedblright{}
\end{quote}

\emph{Th-} becomes \emph{d-}, exactly as it did in \emph{drei}. Take the \emph{th} off \emph{thunder}, put a \emph{d} there, and you have most of the German word already.

This is the second time that rule has paid, and it will keep paying: German \emph{d} answers English \emph{th} right across the vocabulary.

\begin{sounds}
\begin{itemize}
\raggedright
  \item \texttt{vowel-o} — short and open, the \emph{o} of English \emph{on}.
  \item \texttt{r-uvular-german} — the \emph{-er} at the end is barely a consonant; it colours the vowel and stops. \emph{DON-ah}, not \emph{DON-err}.
\end{itemize}
\end{sounds}

\subsection*{Your turn}
Say these aloud:

\begin{itemize}
\raggedright
  \item \textquotedblleft{}der Donner\textquotedblright{}
  \item the rule that predicts it — \textquotedblleft{}th becomes d: three, drei; thunder, Donner\textquotedblright{}
  \item the soft ending — \textquotedblleft{}DON-ah\textquotedblright{}
\end{itemize}

\emph{Twice through:} \textquotedblleft{}der Donner — thunder.\textquotedblright{}

\begin{tcolorbox}[breakable,colback=teal!4,colframe=teal!35!black,title={Before you move on}]
Say \textquotedblleft{}thunder.\textquotedblright{} (\textbf{Der Donner}.) Which English sound does German \emph{d} answer? (\emph{\textbf{Th}}.) Name the other word you have that shows it. (\emph{\textbf{Drei}} / \emph{three}.) Next: the day this noun names, and the god who is both.
\end{tcolorbox}
\section[Donnerstag]{Donnerstag — \textquotedblleft{}Thursday\textquotedblright{}}

\label{lesson:GE-C07-donnerstag}



\begin{quote}
You have the thunder and you have the pattern. This one you can build.
\end{quote}

\subsection*{You'll want to know: Donnerstag}
\begin{quote}
\textbf{Donnerstag} — \textquotedblleft{}Thursday.\textquotedblright{}
\end{quote}

\emph{Donner} plus \emph{-s-} plus \emph{Tag}. The little \emph{s} in the middle is a joining sound German drops between the halves of a long compound; it carries no meaning of its own.

\begin{cousinweb}
The thunder-god was \textbf{Donar} in the German-speaking south and \textbf{Thor} in the Norse north — one god, two names, and \emph{Donner} is his name still doing duty as an ordinary weather word.

English named the day for the Norse form: \textbf{Thurs}day is \emph{Thor's day}.

He stood in for \textbf{Jupiter}, who was the Roman sky-and-thunder king. So this one slot has been filled three times over — Jupiter, Donar, Thor — and every version of it is the god who throws the lightning.

Say \emph{Donnerstag} and \emph{Thursday} back to back and the shift you already own does the rest: German \emph{d} where English has \emph{th}.
\end{cousinweb}

\subsection*{Your turn}
Say these aloud:

\begin{itemize}
\raggedright
  \item \textquotedblleft{}Donnerstag\textquotedblright{}
  \item its pieces — \textquotedblleft{}Donner, s, Tag\textquotedblright{}
  \item the three names for one god — \textquotedblleft{}Donar, Thor, Jupiter\textquotedblright{}
\end{itemize}

\emph{Twice through:} \textquotedblleft{}Donnerstag — Thursday.\textquotedblright{}

\begin{tcolorbox}[breakable,colback=teal!4,colframe=teal!35!black,title={Before you move on}]
Say \textquotedblleft{}Thursday.\textquotedblright{} (\textbf{Donnerstag}.) Which god is in it? (\textbf{Donar} — \textbf{Thor} in the north.) Which Roman god did he replace? (\textbf{Jupiter}.) What is the \emph{s} in the middle doing? (\textbf{Joining the two halves}; it means nothing.) Next: the last of the weekdays, and a goddess.
\end{tcolorbox}
\section[Freitag]{Freitag — \textquotedblleft{}Friday\textquotedblright{}}

\label{lesson:GE-C07-freitag}



\begin{quote}
The last weekday, and the \emph{ei} rule you learned with \emph{eins} is all you need for the sound.
\end{quote}

\subsection*{You'll want to know: Freitag}
\begin{quote}
\textbf{Freitag} — \textquotedblleft{}Friday.\textquotedblright{}
\end{quote}

\emph{Ei} is the English word \emph{eye}, always — so the first syllable rhymes with \emph{fry}, never with \emph{free}. \emph{FRY-tahk}.

That one rule has now told you how to say \emph{eins}, \emph{zwei}, \emph{drei} and this. It does not have exceptions to remember.

\begin{cousinweb}
The goddess is \textbf{Frija} — \textbf{Frigg} in the Norse — and she stood in for \textbf{Venus}. English named the day for her too: \textbf{Fri}day.

This is the one day of the week named for a \textbf{goddess}, and it is a goddess in both systems. The Germanic swap did not merely fill the slot; it matched the god in it. Whoever made these substitutions was paying attention to who the Roman deity was, not only to which day was which.
\end{cousinweb}

\subsection*{Your turn}
Say these aloud:

\begin{itemize}
\raggedright
  \item \textquotedblleft{}Freitag\textquotedblright{} — \emph{FRY-tahk}
  \item the goddess — \textquotedblleft{}Frija, Frigg\textquotedblright{}
  \item who she replaced — \textquotedblleft{}Venus\textquotedblright{}
  \item the four you own — \textquotedblleft{}Montag, Dienstag, Donnerstag, Freitag\textquotedblright{}
\end{itemize}

\emph{Twice through:} \textquotedblleft{}Freitag — Friday.\textquotedblright{}

\begin{tcolorbox}[breakable,colback=teal!4,colframe=teal!35!black,title={Before you move on}]
Say \textquotedblleft{}Friday.\textquotedblright{} (\textbf{Freitag}.) How does \emph{ei} sound? (\textbf{Like English \emph{eye}}.) Which goddess is in it? (\textbf{Frija}, or \emph{Frigg}.) Which Roman goddess did she replace? (\textbf{Venus}.) Next: the day that broke the pattern.
\end{tcolorbox}
\section[Mittwoch]{Mittwoch — the day with no god}

\label{lesson:GE-C07-mittwoch}



\begin{quote}
Four days in and the shape has held every time. Here is the first one that breaks it, and you can hear the break before you know the reason: this word does not end in \emph{Tag} at all.
\end{quote}

\subsection*{You'll want to know: Mittwoch}
\begin{quote}
\textbf{Mittwoch} — \textquotedblleft{}Wednesday,\textquotedblright{} and literally \textquotedblleft{}\textbf{mid-week}.\textquotedblright{}
\end{quote}

No god, no \emph{Tag}. \emph{Mitt-} is the middle, \emph{Woche} is the week, and the day is named for its position rather than for anybody.

English did \textbf{not} do this. \emph{Wednesday} is \textbf{Woden's day} — the chief god, sitting exactly where the pattern says he should. So English kept the Germanic name and German threw it away, which is the opposite of what the rest of this chapter would lead you to expect.

\begin{sounds}
\begin{itemize}
\raggedright
  \item \texttt{ch-ach} — the ending is the raspy back-of-the-throat sound you first met in \emph{Nacht}. Same sound, same spelling, and it closes this word too: \emph{MIT-vokh}.
\end{itemize}

The \emph{w} in the middle is a German \emph{w}, so it is a \emph{v}.
\end{sounds}

\begin{culture}
The medieval Church disliked naming a day of its week for the chief pagan god. So German dropped Wodan and put a plain description in his place, and the replacement stuck.

That is a decision you can date and attribute, sitting inside an everyday word. Most of the etymology in this book is drift — sounds wearing down over centuries with nobody deciding anything. This one is somebody choosing, and the choice is still in the language a thousand years later.

English, further from that authority, kept its god.
\end{culture}

\subsection*{Your turn}
Say these aloud:

\begin{itemize}
\raggedright
  \item \textquotedblleft{}Mittwoch\textquotedblright{} — \emph{MIT-vokh}, with the \emph{Nacht} sound at the end
  \item what it means — \textquotedblleft{}mid-week\textquotedblright{}
  \item the contrast — \textquotedblleft{}Mittwoch, Wednesday: German dropped the god, English kept him\textquotedblright{}
\end{itemize}

\emph{Twice through:} \textquotedblleft{}Mittwoch — mid-week.\textquotedblright{}

\begin{tcolorbox}[breakable,colback=teal!4,colframe=teal!35!black,title={Before you move on}]
Say \textquotedblleft{}Wednesday.\textquotedblright{} (\textbf{Mittwoch}.) What does it literally mean? (\textbf{Mid-week}.) Does it end in \emph{Tag}? (\textbf{No} — the only weekday that does not.) Which god does English keep there? (\textbf{Woden}.) Why did German drop him? (\textbf{The Church declined to name a day for the chief pagan god}.) Next: the five weekdays run together.
\end{tcolorbox}
\section[Practice]{Practice — the working week}

\label{lesson:GE-C07-wochentage-1}



\begin{quote}
No new words. Five days, four gods, and one day where the god was taken out.
\end{quote}

\begin{grammarlens}[title={Grammar Lens: the five, assembled}]
\noindent
\begin{tabularx}{\linewidth}{@{}>{\raggedright\arraybackslash}X>{\raggedright\arraybackslash}X@{}}
\toprule
\textbf{German} & \textbf{who is in it} \\
\midrule
\emph{Montag} & the moon \\
\emph{Dienstag} & Tiw, the war-god \\
\emph{Mittwoch} & nobody — \textquotedblleft{}mid-week\textquotedblright{} \\
\emph{Donnerstag} & Donar, the thunder-god \\
\emph{Freitag} & Frija, the goddess \\
\bottomrule
\end{tabularx}

Four of the five carry a god and end in \emph{Tag}. The one that does neither is the one the Church rewrote.
\end{grammarlens}

\subsection*{Your turn}
Say these aloud:

\begin{itemize}
\raggedright
  \item the five in order — \textquotedblleft{}Montag, Dienstag, Mittwoch, Donnerstag, Freitag\textquotedblright{}
  \item them against the English — \textquotedblleft{}Monday, Tuesday, Wednesday, Thursday, Friday\textquotedblright{}
  \item the gods — \textquotedblleft{}the moon, Tiw, nobody, Donar, Frija\textquotedblright{}
  \item the two nouns you own underneath — \textquotedblleft{}der Mond, der Donner\textquotedblright{}
\end{itemize}

\emph{Twice through:} \textquotedblleft{}Montag, Dienstag, Mittwoch, Donnerstag, Freitag.\textquotedblright{}

\begin{tcolorbox}[breakable,colback=teal!4,colframe=teal!35!black,title={Before you move on}]
Run the five without prompting. Which one has no god in it? (\emph{\textbf{Mittwoch}}.) Which has a goddess? (\emph{\textbf{Freitag}} — Frija.) Which two are built on nouns you can say on their own? (\emph{\textbf{Montag}} \textbf{from} \emph{Mond}, \emph{\textbf{Donnerstag}} \textbf{from} \emph{Donner}.) Next: the weekend, where the pattern breaks a second time.
\end{tcolorbox}
